\section*{Bab \ref{ch:b1:spherical-astronomy} --- Astronomi Bola}

\begin{solution}{exo:b1:spherical-astronomy:1}
\omterm{def:b1:spherical-astronomy:horizontal}{Tinggi kutub} sama dengan \omterm{def:b1:spherical-astronomy:geographic}{lintang}: \ang{52}. Sebuah bintang tidak pernah terbenam jika
$\delta > \ang{90} - \ang{52} = \ang{38}$, dan tidak pernah terbit jika
$\delta < \ang{-38}$.
\end{solution}

\begin{solution}{exo:b1:spherical-astronomy:2}
$\qty{6.75}{h}\times\ang{15} = \ang{101.25}$. $H = \mathrm{LST} - \alpha =
\qty{10}{h} - \qty{6}{h}\,\qty{45}{min} = \qty{3}{h}\,\qty{15}{min}$, yaitu
\ang{48.75}. \omterm{def:b1:spherical-astronomy:equatorial}{Sudut jamnya} positif: bintang berada di barat meridian, yang telah ia lintasi tiga jam seperempat \omterm{def:b1:spherical-astronomy:sidereal-time}{waktu sideris} yang lalu.
\end{solution}

\begin{solution}{exo:b1:spherical-astronomy:3}
$h = \ang{90} - \varphi + \delta_\odot$: \ang{73.44} pada solstis Juni, \ang{50} pada ekuinoks, \ang{26.56} pada solstis Desember.
\end{solution}

\begin{solution}{exo:b1:spherical-astronomy:4}
1900: habis dibagi 100 tetapi tidak habis dibagi 400, bukan tahun kabisat. 2000: habis dibagi 400, tahun kabisat. 2100: bukan tahun kabisat. Tahun Julius melebihi \omterm{def:b1:spherical-astronomy:tropical-year}{tahun tropis} sebesar $365.25 - 365.2422 = 0.0078$ hari, yaitu \num{0.78} hari per abad; tahun Gregorius sebesar $0.0003$ hari, yaitu \num{0.03} hari per abad.
\end{solution}

\begin{solution}{exo:b1:spherical-astronomy:5}
$\cos\Delta = \sin^2\ang{45} + \cos^2\ang{45}\cos\ang{90} = 0.5$, sehingga
$\Delta = \ang{60}$ dan jaraknya $6371\times\pi/3 = \qty{6672}{km}$.
Sepanjang paralel, \omterm{def:b1:spherical-astronomy:spherical-triangle}{lingkaran kecil} berjari-jari $R\cos\ang{45}$, jaraknya adalah
$6371\times\cos\ang{45}\times\pi/2 = \qty{7076}{km}$. \omterm{def:b1:spherical-astronomy:spherical-triangle}{Lingkaran besar} lebih
pendek sekitar \qty{404}{km}.
\end{solution}

\begin{solution}{exo:b1:spherical-astronomy:6}
$\sin h = \sin\ang{48}\sin\ang{20} + \cos\ang{48}\cos\ang{20}\cos\ang{30}
= 0.2542 + 0.5446 = 0.7987$, $h = \ang{53.0}$.
$\cos\mathrm{A} = (\sin\ang{20} - \sin\ang{48}\sin\ang{53.0})/(\cos\ang{48}
\cos\ang{53.0}) = (0.3420 - 0.5935)/0.4027 = -0.6246$, $\arccos = \ang{128.7}$;
bintang berada di barat meridian ($H > 0$), sehingga $\mathrm{A} = \ang{360} -
\ang{128.7} = \ang{231.3}$: ke arah barat daya, \ang{53} di atas cakrawala.
\end{solution}

\begin{solution}{exo:b1:spherical-astronomy:7}
$\cos H_0 = -\tan\ang{60}\tan\ang{23.44} = -1.7321\times0.4335 = -0.7510$,
$H_0 = \ang{138.67}$, panjang hari $2\times138.67/15 = \qty{18.49}{h}$, sekitar
\qty{18}{h}~\qty{29}{min}, seperti yang ditunjukkan kurva \ang{60} LU pada solstis Juni.
\end{solution}

\begin{solution}{exo:b1:spherical-astronomy:8}
Lintasan itu bergeser lebih awal sebesar \qty{3.93}{min} setiap hari, sehingga dalam 30 hari sebesar $30\times3.93 =
\qty{118}{min}$: bintang itu melintas sekitar pukul 20.02. Agar melintas pada tengah hari, ia harus bergeser \qty{10}{h} = \qty{600}{min} lebih awal, yang memerlukan
$600/3.93 \approx 153$ hari.
\end{solution}

\begin{solution}{exo:b1:spherical-astronomy:9}
(a) Pada ekuinoks \omterm{def:b1:spherical-astronomy:horizontal}{tinggi tengah hari} adalah \ang{45}: $I = 1361\times\sin\ang{45}
= \qty{962}{W.m^{-2}}$. (b) Pada solstis Juni $h = \ang{68.44}$ dan
$(a/r)^2 = (149.598/152.100)^2 = 0.9674$:
$I = 1361\times0.9674\times\sin\ang{68.44} = \qty{1224}{W.m^{-2}}$. Matahari yang
lebih tinggi lebih dari cukup mengimbangi jarak yang lebih besar.
\end{solution}

\begin{solution}{exo:b1:spherical-astronomy:10}
$\sin\delta_\odot = \sin\ang{23.44}\sin\ang{60} = 0.3445$,
$\delta_\odot = \ang{20.15}$. $\tan\alpha_\odot = \cos\ang{23.44}\tan\ang{60}
= 1.5890$, $\alpha_\odot = \ang{57.82}$ (kuadran yang sama dengan $\lambda_\odot$).
\omterm{def:b1:spherical-astronomy:equatorial}{Asensio rekta} tertinggal sebesar \ang{2.18}, yaitu $2.18\times4 \approx
\qty{8.7}{min}$ waktu. Matahari yang bergerak seragam sepanjang \omterm{def:b1:spherical-astronomy:ecliptic}{ekliptika} tetap bergerak tidak merata dalam \omterm{def:b1:spherical-astronomy:equatorial}{asensio rekta}, besaran yang menentukan \omterm{def:b1:spherical-astronomy:equatorial}{sudut jam} dan karenanya jam matahari: di belakang Matahari rata-rata antara ekuinoks dan solstis berikutnya, di depannya antara solstis dan ekuinoks berikutnya. Inilah suku kemiringan dalam \omterm{def:b1:spherical-astronomy:solar-time}{perata waktu}.
\end{solution}

\begin{solution}{exo:b1:spherical-astronomy:11}
Saat terbit $h = 0$, sehingga \cref{prop:b1:spherical-astronomy:transform} memberikan
$\sin\delta = \cos\varphi\cos\mathrm{A}_0$; pada solstis Juni
$\delta = \varepsilon$, sehingga $\cos\mathrm{A}_0 = \sin\varepsilon/\cos\varphi$.
Untuk $\mathrm{A}_0 = \ang{45}$: $\cos\varphi = \sin\ang{23.44}/\cos\ang{45} =
0.5626$, $\varphi = \ang{55.8}$ LU. Di atas $\ang{90} - \varepsilon =
\ang{66.56}$, $\sin\varepsilon/\cos\varphi > 1$: tidak ada \omterm{def:b1:spherical-astronomy:horizontal}{azimut} terbit, karena
Matahari tidak terbenam pada hari itu (matahari tengah malam).
\end{solution}

\begin{solution}{exo:b1:spherical-astronomy:12}
(a) Baji dengan sudut $\theta$ adalah pecahan $\theta/2\pi$ dari bola:
$\frac{\theta}{2\pi}\times4\pi R^2 = 2\theta R^2$. (b) Ketiga pasang baji
memiliki luas total $2\times(2A + 2B + 2C)R^2$; mereka menutupi bola satu kali,
kecuali segitiga dan segitiga seberangnya, yang masing-masing tertutup tiga kali, yaitu
dua kali lebih banyak daripada yang diperlukan: $4(A + B + C)R^2 = 4\pi R^2 + 4S$, dengan $S$ luas
segitiga. Maka $S = R^2(A + B + C - \pi)$. (c) Tiga sudut siku-siku:
$S = R^2(3\pi/2 - \pi) = \pi R^2/2$, seperdelapan bola (sebuah oktan).
Segitiga pada gambar: $89.9 + 64.9 + 67.4 - 180 = \ang{42.2}$, yaitu
\num{0.736} radian: $S = 0.736\,R^2$, \qty{5.9}{\percent} dari bola.
\end{solution}

\begin{solution}{pb:b1:spherical-astronomy:1}
\textbf{1.} Kolatitud $\ang{90} - (\ang{-24.63}) = \ang{114.63}$ dan
$\ang{90} - \ang{51.48} = \ang{38.52}$; selisih \omterm{def:b1:spherical-astronomy:geographic}{bujur} \ang{70.40}.

\textbf{2.} $\cos\Delta = \cos\ang{114.63}\cos\ang{38.52} +
\sin\ang{114.63}\sin\ang{38.52}\cos\ang{70.40} = -0.3260 + 0.1899 =
-0.1362$.

\textbf{3.} $\Delta = \ang{97.83}$; jarak $6371\times1.7074 =
\qty{10878}{km}$, sekitar \qty{10900}{km}.

\textbf{4.}
\[
  \cos\mathrm{Az} = \frac{\cos\ang{38.52} - \cos\ang{114.63}\cos\ang{97.83}}
  {\sin\ang{114.63}\sin\ang{97.83}} = \frac{0.7823 - 0.0568}{0.9006} = 0.8056 ,
\]
$\mathrm{Az} = \ang{36.3}$, ke arah timur laut.

\textbf{5.} $\sin\mathrm{Az} = \sin\ang{70.40}\sin\ang{38.52}/\sin\ang{97.83}
= 0.5922$, yang memberikan \ang{36.3} atau \ang{143.7}; kosinus positif pada
soal 4 memilih \ang{36.3}.

\textbf{6.} Tidak. Arah yang tetap memotong setiap meridian dengan sudut yang sama dan
melingkar menuju kutub; arah \omterm{def:b1:spherical-astronomy:spherical-triangle}{lingkaran besar} berubah terus-menerus
(hanya \ang{36.3} pada awalnya). Pesawat itu akan menempuh jalan yang lebih panjang
dan tiba di tempat lain kecuali ia mengoreksi arahnya.

\textbf{7.} \omterm{def:b1:spherical-astronomy:horizontal}{Tinggi} \omterm{def:b1:spherical-astronomy:celestial-sphere}{kutub langit} selatan sama dengan besar \omterm{def:b1:spherical-astronomy:geographic}{lintangnya}: \ang{24.63}.

\textbf{8.} $h = \ang{90} - |\ang{-24.63} - (\ang{-16.72})| = \ang{82.09}$;
$\delta > \varphi$, sehingga Sirius melintas \ang{7.91} di utara \omterm{def:b1:spherical-astronomy:celestial-sphere}{zenit}.

\textbf{9.} Sebuah bintang tidak pernah terbenam jika $\delta < -(\ang{90} - \ang{24.63}) =
\ang{-65.37}$. Canopus, pada \ang{-52.70}, tidak sirkumpolar: ia terbenam.

\textbf{10.} $\ang{90} - |\ang{-24.63} + \ang{52.70}| = \ang{61.93}$, di selatan
\omterm{def:b1:spherical-astronomy:celestial-sphere}{zenit}.

\textbf{11.} $\cos H_0 = -\tan(\ang{-24.63})\tan(\ang{-16.72}) =
-(0.4584)(0.3004) = -0.1377$; $H_0 = \ang{97.92}$. Sirius tetap di atas cakrawala
selama $2\times97.92/15 = \qty{13.06}{h}$ \omterm{def:b1:spherical-astronomy:sidereal-time}{waktu sideris}, sekitar \qty{13}{h}.

\textbf{12.} $\cos\mathrm{A}_0 = \sin(\ang{-16.72})/\cos\ang{24.63} =
-0.3165$: Sirius terbit pada \omterm{def:b1:spherical-astronomy:horizontal}{azimut} \ang{108.5}, di selatan timur.

\textbf{13.} Ketika $H = 0$, $\mathrm{LST} = \alpha = \qty{6}{h}~\qty{45}{min}~\qty{9}{s}$.

\textbf{14.} $\qty{86400}{s}\times365.2422/366.2422 = \qty{86164.09}{s} =
\qty{23}{h}~\qty{56}{min}~\qty{4}{s}$; lintasan bergeser lebih awal
\qty{235.9}{s}, yaitu \qty{3}{min}~\qty{56}{s}, setiap malam.

\textbf{15.} Kita memerlukan $12 + n\times24/365.2422 = 6.7525 + 24$, sehingga
$n = 18.7525/0.065710 = 285.4$ hari. 285 hari setelah 20 Maret jatuh pada
30 Desember: Sirius melintasi meridian pada tengah malam \omterm{def:b1:spherical-astronomy:solar-time}{waktu matahari rata-rata} pada malam
30 menuju 31 Desember 2026.

\textbf{16.} Paranal berada \ang{70.40} di barat Greenwich: \omterm{def:b1:spherical-astronomy:solar-time}{waktu matahari rata-rata}
tertinggal sebesar $70.40/15 = \qty{4.693}{h}$, yaitu \qty{4}{h}~\qty{41.6}{min}.

\textbf{17.} \omterm{def:b1:spherical-astronomy:solar-time}{Waktu matahari semu} = \omterm{def:b1:spherical-astronomy:solar-time}{waktu matahari rata-rata} + \omterm{def:b1:spherical-astronomy:solar-time}{perata waktu} =
$\qty{0}{h} - \qty{3}{min}$: jam matahari menunjukkan sekitar 23.57.

\textbf{18.} $7\times\qty{3.93}{min} \approx \qty{27.5}{min}$ lebih awal.

\textbf{19.} $+\ang{23.44}$ pada solstis Juni, $-\ang{23.44}$ pada
solstis Desember.

\textbf{20.} Juni: $\ang{90} - |\ang{-24.63} - \ang{23.44}| = \ang{41.93}$.
Desember: $\ang{90} - |\ang{-24.63} + \ang{23.44}| = \ang{88.81}$; karena
$\delta_\odot = \ang{-23.44} > \varphi$, Matahari melintas \ang{1.19} di utara
\omterm{def:b1:spherical-astronomy:celestial-sphere}{zenit}.

\textbf{21.} $I = 1361\times(1/0.9837)^2\times\sin\ang{88.81} = 1361\times
1.0334\times0.9998 = \qty{1406}{W.m^{-2}}$.

\textbf{22.} Pada bulan Juni, musim dingin di belahan selatan. $\cos H_0 = -\tan(\ang{-24.63})
\tan\ang{23.44} = 0.4584\times0.4335 = 0.1988$: Matahari terbenam pada \omterm{def:b1:spherical-astronomy:equatorial}{sudut jam}
$H_0 = \ang{78.53}$.

\textbf{23.} $\cos\mathrm{A}_0 = \sin\ang{23.44}/\cos\ang{24.63} = 0.4376$:
$\mathrm{A}_0 = \ang{64.0}$, di utara timur.

\textbf{24.} Panjang hari $2\times78.53/15 = \qty{10.47}{h}$; malam terpanjang
dalam setahun di Paranal berlangsung $24 - 10.47 = \qty{13.53}{h}$, yaitu
\textbf{\qty[text-series-to-math=true]{13}{h}~\qty[text-series-to-math=true]{32}{min}}.
\end{solution}
